% =====================================================================
%  CARNET D'ENTRAINEMENT — FICHE 6 — THÈME 1
%  Niveau DIFFICILE — Exercices
% =====================================================================
\documentclass[11pt,a4paper]{article}
\usepackage[utf8]{inputenc}
\usepackage[T1]{fontenc}
\usepackage{lmodern}
\usepackage[french]{babel}
\usepackage{amsmath,amssymb,mathtools}
\usepackage{icomma}
\usepackage[version=4]{mhchem}
\usepackage{siunitx}
\sisetup{output-decimal-marker={,}}
\usepackage{xcolor}
\usepackage{colortbl}
\usepackage{tikz}
\usepackage[most]{tcolorbox}
\usepackage{enumitem}
\usepackage{array}
\usepackage{booktabs}
\usepackage{fancyhdr}
\usepackage[hidelinks]{hyperref}
\usetikzlibrary{arrows.meta,positioning,calc,shapes.geometric}
\usepackage[a4paper,margin=2.1cm,headheight=26pt,headsep=8pt]{geometry}

\definecolor{primary}{RGB}{146,95,160}
\definecolor{primarydark}{RGB}{92,52,104}
\definecolor{atomCl}{RGB}{70,170,80}
\definecolor{atomNa}{RGB}{150,90,210}
\definecolor{bondcol}{RGB}{60,60,70}
\definecolor{piecol}{RGB}{198,93,9}

\tcbset{boxbase/.style={enhanced,breakable,boxrule=0.9pt,arc=2.5pt,
   left=3mm,right=3mm,top=2.2mm,bottom=2.2mm,
   fonttitle=\bfseries,coltitle=white,toptitle=0.6mm,bottomtitle=0.6mm}}
\newtcolorbox[auto counter]{exo}[1][]{boxbase,colback=primary!5,colframe=primary,
   title={\textsc{Exercice}~\thetcbcounter\ifx\relax#1\relax\relax\else\textmd{ —\itshape\ #1}\fi}}

\pagestyle{fancy}\fancyhf{}
\fancyhead[L]{\small\textcolor{primarydark}{\textbf{Fiche 6 · Thème 1} — Types de liaisons et octet}}
\fancyhead[R]{\small\textcolor{primarydark}{\textsc{Difficile}}}
\fancyfoot[C]{\small\thepage}
\renewcommand{\headrulewidth}{0.8pt}
\renewcommand{\headrule}{\hbox to\headwidth{\color{primary}\leaders\hrule height \headrulewidth\hfill}}
\setlength{\parindent}{0pt}\setlength{\parskip}{3pt}

\begin{document}
\begin{center}
{\Large\bfseries\color{primarydark}Thème 1 — Niveau Difficile}\\[2pt]
{\small\itshape Les questions sont décomposées en étapes (a, b, c\ldots) qui construisent le raisonnement.}
\end{center}
\medskip

\begin{exo}[la formation du sel de table, pas à pas]
Le chlorure de sodium \ce{NaCl} (le sel de table) est un solide ionique. On étudie la formation
d'une paire \ce{Na}\,--\,\ce{Cl}.
\begin{enumerate}[label=\alph*),leftmargin=2em,itemsep=2pt]
  \item Donner le nombre d'électrons de valence de \ce{Na} et de \ce{Cl}.
  \item Pour atteindre une couche externe complète, lequel des deux atomes a intérêt à
        \textbf{céder} un électron, et lequel à le \textbf{capter} ? Écrire les deux ions formés.
  \item Vérifier que chaque ion possède alors une configuration stable : préciser, pour chacun,
        s'il atteint un octet ou un duet, et le gaz rare correspondant.
  \item Quel est le \textbf{type de liaison} qui assure la cohésion du solide, et par quelle force
        physique ?
\end{enumerate}
\end{exo}

\begin{exo}[l'ion ammonium et la liaison dative]
On s'intéresse à la réaction \ce{NH3 + H+ -> NH4+}.
\begin{enumerate}[label=\alph*),leftmargin=2em,itemsep=2pt]
  \item Dans \ce{NH3}, l'azote (5 électrons de valence) forme 3 liaisons \ce{N-H}. En déduire
        combien de \textbf{doublets non liants} il porte.
  \item Expliquer, en utilisant les mots \emph{donneur} et \emph{accepteur}, comment se forme la
        quatrième liaison \ce{N-H}.
  \item Pourquoi, une fois l'ion \ce{NH4+} formé, les quatre liaisons \ce{N-H} sont-elles
        strictement \textbf{identiques} ?
  \item L'azote respecte-t-il l'octet dans \ce{NH4+} ? Compter les électrons qui l'entourent.
\end{enumerate}
\end{exo}

\begin{exo}[trois atomes centraux, trois comportements face à l'octet]
On compare \ce{BF3}, \ce{CF4} et \ce{SF6}. Dans chacun, l'atome central ne forme que des liaisons
simples et ne porte aucun doublet libre.
\begin{enumerate}[label=\alph*),leftmargin=2em,itemsep=2pt]
  \item Pour chaque molécule, compter le nombre d'électrons autour de l'atome central.
  \item Classer chaque cas : octet \textbf{incomplet}, \textbf{respecté}, ou \textbf{étendu}.
  \item Justifier pourquoi l'écart à l'octet est \emph{permis} pour \ce{BF3} d'une part, pour
        \ce{SF6} d'autre part.
  \item Une molécule hypothétique « \ce{OF6} » (6 liaisons autour de l'oxygène) pourrait-elle
        exister ? Justifier.
\end{enumerate}
\end{exo}

\begin{exo}[le cristal ionique : vrai ou faux]
Le chlorure de sodium existe à l'état solide sous forme d'un \textbf{cristal ionique}. Pour chaque
affirmation, répondre par \textbf{vrai} ou \textbf{faux} en justifiant.
\begin{enumerate}[label=\alph*),leftmargin=2em,itemsep=2pt]
  \item « La cohésion d'un cristal ionique est due à des interactions gravitationnelles. »
  \item « Un cristal ionique possède une structure ordonnée, et cet ordre est détruit lors de la
        dissolution dans l'eau. »
  \item « La solution obtenue après dissolution conduit le courant électrique. »
\end{enumerate}
\end{exo}

\end{document}
